\chapter{Integration Techniques}

\begin{center}\itshape
I had learned to do integrals by various methods shown in a book that my high school physics teacher\\
gave me. One thing I never did learn was contour integration. [\dots] I had a different box of tools.\\
\upshape --- Richard Feynman, \emph{Surely You're Joking, Mr.~Feynman!}
\end{center}

Feynman's ``different box of tools'' was mostly one tool: put a parameter $t$ into the integral, differentiate with
respect to $t$ under the integral sign, and solve the differential equation that comes out. With it he could compute
integrals that defeated his colleagues at MIT and Princeton. The trick is two centuries older (Leibniz used it, and so
did Euler), and it always needs a justification, since swapping a derivative with an integral can fail.

This chapter supplies the justification and then uses it. The first two sections prove when limits and derivatives
pass under an integral sign, first on a closed bounded interval, where uniform continuity does everything, then on
unbounded intervals, where an integrable bound independent of the parameter is needed; they also prove Fubini's
theorem for continuous functions on rectangles. The rest of the chapter is a tour of the toolbox: Feynman's trick,
Frullani's integrals, symmetry, and the most famous semi-convergent integral, Dirichlet's
$\int_0^\infty\frac{\sin x}x\dd x=\frac\pi2$. Complex analysis, the other box of tools, comes in Chapter~9.

\begin{think}
For $t>0$, the substitution $u=tx$ shows that $\int_0^\infty\frac{\sin tx}x\dd x$ does not depend on $t$. Differentiating
under the integral sign gives $\int_0^\infty\cos tx\dd x$. What has gone wrong?
\end{think}

\section{Parameters on a compact interval}

Tome~0 integrated continuous functions on closed bounded intervals. With a parameter, the key is that continuity on a
closed bounded rectangle is automatically uniform.

\begin{lemma}[Heine]\label{lem:heine}
A continuous function $f$ on a closed bounded rectangle $K=[a,b]\times[c,d]$ is uniformly continuous: for every
$\varepsilon>0$ there is $\delta>0$ such that $|f(p)-f(q)|\le\varepsilon$ whenever $p,q\in K$ and $|p-q|\le\delta$.
\end{lemma}
\begin{proof}
Otherwise there are $\varepsilon>0$ and points $p_n,q_n\in K$ with $|p_n-q_n|\le\frac1n$ and $|f(p_n)-f(q_n)|>\varepsilon$. Applying
Bolzano--Weierstrass (Tome~0, Chapter~5) to the first coordinates, then to the second coordinates along the extracted
subsequence, we get a subsequence $p_{n_k}\to p\in K$, and then $q_{n_k}\to p$ too. By continuity
$f(p_{n_k})-f(q_{n_k})\to f(p)-f(p)=0$, a contradiction.
\end{proof}

\begin{theorem}\label{thm:paramcompact}
Let $J$ be an interval and $f$ continuous on $[a,b]\times J$. Then $F(t)=\int_a^bf(x,t)\dd x$ is continuous on $J$. If
moreover $\partial_tf$ exists and is continuous on $[a,b]\times J$, then $F$ is of class $C^1$ and
\[
F'(t)=\int_a^b\partial_tf(x,t)\dd x.
\]
\end{theorem}
\begin{proof}
Let $t_0\in J$ and let $[c,d]\subseteq J$ be a closed bounded interval that is a neighbourhood of $t_0$ in $J$. Given $\varepsilon>0$,
Lemma~\ref{lem:heine} on $[a,b]\times[c,d]$ gives $\delta$ with $|f(x,t)-f(x,t_0)|\le\varepsilon$ for $|t-t_0|\le\delta$, $t\in[c,d]$, so
$|F(t)-F(t_0)|\le(b-a)\varepsilon$. For the derivative, let $q(x,t)=\frac{f(x,t)-f(x,t_0)}{t-t_0}$. By the mean value theorem, for each $x$
there is $\tau$ between $t_0$ and $t$ with
\[
q(x,t)-\partial_tf(x,t_0)=\partial_tf(x,\tau)-\partial_tf(x,t_0),
\]
which is at most $\varepsilon$ in absolute value when $|t-t_0|\le\delta$, now with $\delta$ given by the uniform continuity of $\partial_tf$
on $[a,b]\times[c,d]$. Integrating, $\bigl|\frac{F(t)-F(t_0)}{t-t_0}-\int_a^b\partial_tf(x,t_0)\dd x\bigr|\le(b-a)\varepsilon$. So $F'(t_0)$ exists and is the
integral of $\partial_tf(\cdot,t_0)$, which is continuous in $t_0$ by the first part.
\end{proof}

\begin{theorem}[Fubini on a rectangle]\label{thm:fubinirect}
If $f$ is continuous on $[a,b]\times[c,d]$, then
\[
\int_a^b\Bigl(\int_c^df(x,y)\dd y\Bigr)\dd x=\int_c^d\Bigl(\int_a^bf(x,y)\dd x\Bigr)\dd y.
\]
\end{theorem}
\begin{proof}
For $s\in[a,b]$ let
\[
\begin{gathered}\varphi(s)=\int_a^s\Bigl(\int_c^df(x,y)\dd y\Bigr)\dd x,\qq\psi(s)=\int_c^d\Bigl(\int_a^sf(x,y)\dd x\Bigr)\dd y.\end{gathered}
\]
The inner
integral in $\varphi$ is a continuous function of $x$ by Theorem~\ref{thm:paramcompact}, so $\varphi'(s)=\int_c^df(s,y)\dd y$ by the
fundamental theorem of calculus. In $\psi$, the function $g(s,y)=\int_a^sf(x,y)\dd x$ is continuous on $[a,b]\times[c,d]$ (it is
$\int_a^bf(x,y)\mathbf 1_{x\le s}\dd x$; directly, $|g(s,y)-g(s',y')|\le|s-s'|\max|f|+(b-a)\max_x|f(x,y)-f(x,y')|$) and has the
continuous partial derivative $\partial_sg(s,y)=f(s,y)$. By Theorem~\ref{thm:paramcompact}, $\psi'(s)=\int_c^df(s,y)\dd y=\varphi'(s)$.
Since $\varphi(a)=\psi(a)=0$, $\varphi=\psi$, and $s=b$ gives the theorem.
\end{proof}

\section{Improper integrals with a parameter}

Let $I$ be an interval of any kind, with endpoints $\alpha<\beta$ in $[-\infty,+\infty]$, and $g$ continuous on $I$. The
\emph{improper integral} $\int_Ig$ converges if $\int_u^vg$ has a limit as $u\to\alpha$ and $v\to\beta$ independently (with $u$, $v$ in
$I$); by the Cauchy criterion this holds if $\int_I|g|<\infty$, and then $\bigl|\int_Ig\bigr|\le\int_I|g|$: we say $g$ is
\emph{integrable} on $I$. A non-negative continuous $g$ is integrable if and only if the integrals over closed bounded
subintervals are bounded, and a continuous $f$ with $|f|\le g$, $g$ integrable, is integrable (comparison). The main
example: $x^{-s}$ is integrable on $[1,\infty[$ if and only if $s>1$, and on $]0,1]$ if and only if $s<1$.

\begin{theorem}[domination]\label{thm:domination}
Let $I$, $J$ be intervals and $f$ continuous on $I\times J$, with $|f(x,t)|\le g(x)$ for all $(x,t)$, where $g$ is continuous and
integrable on $I$. Then $F(t)=\int_If(x,t)\dd x$ is continuous on $J$. If moreover $\partial_tf$ is continuous on $I\times J$ and
$|\partial_tf(x,t)|\le h(x)$ with $h$ continuous and integrable on $I$, then $F$ is of class $C^1$ and $F'(t)=\int_I\partial_tf(x,t)\dd x$.
\end{theorem}
\begin{proof}
Given $\varepsilon>0$, choose a closed bounded $[u,v]\subseteq I$ with $\int_{I\setminus[u,v]}g\le\varepsilon$ and $\int_{I\setminus[u,v]}h\le\varepsilon$. Then
for $t,t_0\in J$, $|F(t)-F(t_0)|\le2\varepsilon+\bigl|\int_u^v(f(x,t)-f(x,t_0))\dd x\bigr|$, and the last term is at most $\varepsilon$ for $t$ close to
$t_0$ by Theorem~\ref{thm:paramcompact} on $[u,v]$. For the derivative, the difference quotient minus $\int_I\partial_tf(x,t_0)\dd x$
is, by the mean value theorem as before, at most $2\varepsilon$ outside $[u,v]$ (each term is bounded by $h$) plus a term over $[u,v]$
that tends to $0$ by Theorem~\ref{thm:paramcompact}. Continuity of $F'$ follows from the first part applied to $\partial_tf$.
\end{proof}

The same proof gives the version for sequences that is used throughout this tome.

\begin{corollary}[dominated convergence for continuous functions]\label{cor:dct}
Let $f_n$ and $f$ be continuous on $I$, with $f_n\to f$ uniformly on every closed bounded subinterval, and $|f_n|\le g$ with $g$
continuous and integrable on $I$. Then $\int_If_n\to\int_If$.
\end{corollary}
\begin{proof}
$|f|\le g$ too, and $\bigl|\int_I(f_n-f)\bigr|\le2\int_{I\setminus[u,v]}g+(v-u)\sup_{[u,v]}|f_n-f|$; choose $[u,v]$ first, then $n$.
\end{proof}

\begin{remark}
Tome~3, Chapter~5, proved the dominated convergence theorem for the Lebesgue integral, which needs only pointwise
convergence and measurable functions; the continuous versions above are all this tome needs, and they are elementary. The
bound $g$ cannot be dropped: Exercise~\ref{ex:swapfail} shows a sequence of integrals that does not converge to the integral of
the limit.
\end{remark}

\section{Feynman's trick}

\begin{example}\label{ex:logint}
For $t\ge0$ let $F(t)=\int_0^1\frac{x^t-1}{\ln x}\dd x$. On $I=\left]0,1\right[$ and $J=[0,\infty[$ the integrand is continuous, with
$\partial_t\frac{x^t-1}{\ln x}=x^t$. For $t\in[0,T]$ the bounds are $\bigl|\frac{x^t-1}{\ln x}\bigr|\le t\le T$, since $1-e^{-u}\le u$ for
$u=t|\ln x|\ge0$, and $|x^t|\le1$. By Theorem~\ref{thm:domination} (on $J=[0,T]$, for every $T$), $F'(t)=\int_0^1x^t\dd x=\frac1{t+1}$. Since
$F(0)=0$, $F(t)=\ln(1+t)$; in particular
\[
\int_0^1\frac{x-1}{\ln x}\dd x=\ln2.
\]
\end{example}

The pattern is always the same: choose where to put the parameter so that the differentiated integrand is elementary;
compute $F'$; find the constant from a value of $t$ where $F$ is obvious, often $t=0$ or $t\to\infty$.

\begin{example}[the Gaussian integral]\label{ex:gaussfeynman}
Let $F(t)=\Bigl(\int_0^te^{-x^2}\dd x\Bigr)^2$ and $G(t)=\int_0^1\frac{e^{-t^2(1+y^2)}}{1+y^2}\dd y$. By Theorem~\ref{thm:paramcompact},
\[
\begin{aligned}G'(t)&=-2t\int_0^1e^{-t^2(1+y^2)}\dd y\eqbr=-2e^{-t^2}\int_0^te^{-u^2}\dd u=-F'(t),\end{aligned}
\]
with $u=ty$. So $F+G$ is constant, equal to $G(0)=\int_0^1\frac{\mathrm dy}{1+y^2}=\frac\pi4$. As $t\to\infty$, $0\le G(t)\le e^{-t^2}\to0$, so
$\bigl(\int_0^\infty e^{-x^2}\dd x\bigr)^2=\frac\pi4$, and $\int_0^\infty e^{-x^2}\dd x=\frac{\sqrt\pi}2$. Tome~0 found it by Wallis integrals, and Tome~3
by polar coordinates; this third proof uses only a parameter.
\end{example}

\section{Frullani integrals and symmetry}

\begin{theorem}[Frullani]\label{thm:frullani}
Let $f$ be continuous on $[0,\infty[$ with a limit $f(\infty)$ at infinity. For $a,b>0$,
\[
\int_0^\infty\frac{f(ax)-f(bx)}x\dd x=\bigl(f(0)-f(\infty)\bigr)\ln\frac ba.
\]
\end{theorem}
\begin{proof}
For $0<\varepsilon<R$, let $I_{\varepsilon,R}=\int_\varepsilon^R\frac{f(ax)-f(bx)}x\dd x$. Substituting $u=ax$ and $u=bx$ in the two halves,
\[
\begin{aligned}I_{\varepsilon,R}&=\int_{a\varepsilon}^{aR}\frac{f(u)}u\dd u-\int_{b\varepsilon}^{bR}\frac{f(u)}u\dd u\\&=\int_{a\varepsilon}^{b\varepsilon}\frac{f(u)}u\dd u-\int_{aR}^{bR}\frac{f(u)}u\dd u.\end{aligned}
\]
(The difference of the two integrals over the common range cancels; if $a>b$, read $\int_p^q$ as $-\int_q^p$.) On
$[a\varepsilon,b\varepsilon]$, $f(u)=f(0)+o(1)$ as $\varepsilon\to0$, and $\int_{a\varepsilon}^{b\varepsilon}\frac{\mathrm du}u=\ln\frac ba$, so the first term tends to
$f(0)\ln\frac ba$. Likewise the second tends to $f(\infty)\ln\frac ba$ as $R\to\infty$.
\end{proof}

With $f(x)=e^{-x}$: $\int_0^\infty\frac{e^{-ax}-e^{-bx}}x\dd x=\ln\frac ba$. With $f(x)=\arctan x$: $\int_0^\infty\frac{\arctan bx-\arctan ax}x\dd x=\frac\pi2\ln\frac ba$.

Symmetry is the cheapest trick of all. The substitution $x\mapsto a+b-x$ preserves $\int_a^b$, and sometimes turns an
integral into a multiple of itself plus something known.

\begin{proposition}\label{prop:logsin}
The integrals $\int_0^{\pi/2}\ln\sin x\dd x$ and $\int_0^{\pi/2}\ln\cos x\dd x$ are both equal to $-\frac\pi2\ln2$.
\end{proposition}
\begin{proof}
The integrals converge, since $\ln\sin x=\ln x+\ln\frac{\sin x}x$ and $\ln x$ is integrable at $0$. The substitution $x\mapsto\frac\pi2-x$
shows they are equal; call their value $L$. Then, using $\sin x\cos x=\frac{\sin2x}2$,
\[
\begin{aligned}2L&=\int_0^{\pi/2}\ln\frac{\sin2x}2\dd x\eqbr=\frac12\int_0^\pi\ln\sin u\dd u-\frac\pi2\ln2\eqbr=L-\frac\pi2\ln2,\end{aligned}
\]
because $\int_0^\pi\ln\sin u\dd u=2L$ by the symmetry $u\mapsto\pi-u$. Hence $L=-\frac\pi2\ln2$.
\end{proof}

\section{Dirichlet's integral}

The function $\frac{\sin x}x$ is not integrable on $[0,\infty[$: over $[k\pi,(k+1)\pi]$ its absolute value has integral at
least $\frac2{(k+1)\pi}$, and the harmonic series diverges. But its integral converges as an improper integral, thanks to the
cancellations between the positive and negative arches (Figure~\ref{fig:si}).

\begin{lemma}\label{lem:secondmean}
Let $u$ be of class $C^1$, positive and non-increasing on $[A,B]$. Then $\bigl|\int_A^Bu(x)\sin x\dd x\bigr|\le2u(A)$.
\end{lemma}
\begin{proof}
Integrate by parts: $\int_A^Bu\sin=\bigl[-u\cos\bigr]_A^B+\int_A^Bu'\cos$, so the absolute value is at most
$u(A)+u(B)+\int_A^B|u'|=u(A)+u(B)+u(A)-u(B)=2u(A)$, since $u'\le0$.
\end{proof}

\begin{theorem}[Dirichlet]\label{thm:dirichletint}
The improper integral $\int_0^\infty\frac{\sin x}x\dd x$ converges, and equals $\frac\pi2$.
\end{theorem}
\begin{proof}
For $t\ge0$ and $A>0$, Lemma~\ref{lem:secondmean} with $u(x)=\frac{e^{-tx}}x$ gives $\bigl|\int_A^B\frac{e^{-tx}\sin x}x\dd x\bigr|\le\frac2A$
for all $B>A$, uniformly in $t\ge0$. With $t=0$, the Cauchy criterion shows that the integral converges (near $0$ the integrand
extends continuously by $1$). Let $F(t)=\int_0^\infty e^{-tx}\frac{\sin x}x\dd x$.

\emph{For $t>0$.} On $J=[t_0,\infty[$ with $t_0>0$, the integrand and its $t$-derivative $-e^{-tx}\sin x$ are bounded by $e^{-t_0x}$,
which is integrable. By Theorem~\ref{thm:domination}, $F'(t)=-\int_0^\infty e^{-tx}\sin x\dd x=-\frac1{1+t^2}$ (integrate by parts twice, or take
the imaginary part of $\int_0^\infty e^{(-t+i)x}\dd x=\frac1{t-i}$). So $F(t)=C-\arctan t$ on $]0,\infty[$. Since $\bigl|\frac{\sin x}x\bigr|\le1$,
$|F(t)|\le\int_0^\infty e^{-tx}\dd x=\frac1t\to0$, and $C=\frac\pi2$.

\emph{At $t=0$.} For $A>0$ and $t>0$, let $D=\bigl|F(t)-\int_0^\infty\frac{\sin x}x\dd x\bigr|$. Then
\[
\begin{aligned}D&\le\int_0^A(1-e^{-tx})\dd x+\frac2A+\frac2A\eqbr\le\frac{tA^2}2+\frac4A,\end{aligned}
\]
using $1-e^{-tx}\le tx$. Letting $t\to0$, then $A\to\infty$, gives $\int_0^\infty\frac{\sin x}x\dd x=\lim_{t\to0}F(t)=\frac\pi2$.
\end{proof}

\begin{figure}[ht]
\centering
\fig{ch04-si}
\caption{The partial integrals $\operatorname{Si}(X)=\int_0^X\frac{\sin x}x\dd x$ oscillate around $\frac\pi2$ with decreasing amplitude. The
first maximum, at $X=\pi$, is $\operatorname{Si}(\pi)\approx1.852$, about $18\%$ above the limit: this overshoot is the Gibbs phenomenon
of Chapter~6.}
\label{fig:si}
\end{figure}

\begin{figure}[ht]
\centering
\fig{ch04-damped}
\caption{The damped integrands $e^{-tx}\frac{\sin x}x$ for $t=0$, $0.3$ and $1$. For $t>0$ the integral is absolutely convergent and can
be differentiated in $t$; the value at $t=0$ is recovered as a limit.}
\label{fig:damped}
\end{figure}

Now the question of the opening box has an answer. For $t>0$, $\int_0^\infty\frac{\sin tx}x\dd x=\frac\pi2$; for $t<0$ it is $-\frac\pi2$,
and for $t=0$ it is $0$. The function of $t$ is a step, not differentiable at $0$, and the formally differentiated integral
$\int_0^\infty\cos tx\dd x$ diverges for every $t$: no integrable bound $h$ exists, and Theorem~\ref{thm:domination} does not apply.

\section{Exercises}

\begin{exercise}[$\star$]\label{ex:swapfail}
Let $f_n(x)=n^2xe^{-nx}$ on $[0,\infty[$. Show that $f_n\to0$ pointwise and $\int_0^\infty f_n=1$ for all $n$. Which hypothesis of
Corollary~\ref{cor:dct} fails?
\end{exercise}

\begin{exercise}[$\star$]\label{ex:gausscos}
Show that $\int_{-\infty}^\infty e^{-x^2}\cos(2ax)\dd x=\sqrt\pi\,e^{-a^2}$ for every real $a$, by differentiating in $a$ and integrating by
parts.
\end{exercise}

\begin{exercise}[$\star$]\label{ex:sin2}
Show that $\int_0^\infty\frac{\sin^2x}{x^2}\dd x=\frac\pi2$, by parts from Theorem~\ref{thm:dirichletint}.
\end{exercise}

\begin{exercise}[$\star\star$]\label{ex:xsinx}
Show that $\int_0^\pi\frac{x\sin x}{1+\cos^2x}\dd x=\frac{\pi^2}4$.
\end{exercise}

\begin{exercise}[$\star\star$]\label{ex:loglog}
Show that $\int_0^1\ln x\ln(1-x)\dd x=2-\frac{\pi^2}6$.
\end{exercise}

\begin{exercise}[$\star\star$]\label{ex:arctanparam}
For $a\ge0$ let
$J(a)=\int_0^\infty\frac{\ln(1+a^2x^2)}{1+x^2}\dd x$. Show that $J(a)=\pi\ln(1+a)$.
\end{exercise}

\section{Hints}

\begin{hint}{ex:swapfail}
Substitute $u=nx$. Where is the mass of $f_n$?
\end{hint}
\begin{hint}{ex:gausscos}
Call the integral $F(a)$. Differentiating under the integral sign gives
\[
F'(a)=-\int2xe^{-x^2}\sin(2ax)\dd x;
\]
integrate by parts to get $F'=-2aF$.
\end{hint}
\begin{hint}{ex:sin2}
$\frac{\mathrm d}{\mathrm dx}\sin^2x=\sin2x$, and $\frac1{x^2}=-\frac{\mathrm d}{\mathrm dx}\frac1x$.
\end{hint}
\begin{hint}{ex:xsinx}
Substitute $x\mapsto\pi-x$ and add; then $u=\cos x$.
\end{hint}
\begin{hint}{ex:loglog}
Expand $\ln(1-x)=-\sum_{n\ge1}\frac{x^n}n$ and integrate term by term, using $\int_0^1x^n\ln x\dd x=-\frac1{(n+1)^2}$; then
$\frac1{n(n+1)^2}=\frac1n-\frac1{n+1}-\frac1{(n+1)^2}$.
\end{hint}
\begin{hint}{ex:arctanparam}
Differentiate under the integral:
\[
J'(a)=\int_0^\infty\frac{2ax^2}{(1+a^2x^2)(1+x^2)}\dd x.
\]
Then split into partial fractions in $x^2$.
\end{hint}

\section{Solutions}

\begin{solution}{ex:swapfail}
\idea The graphs are bumps of fixed area that slide towards $0$ and grow taller: no single integrable function lies above all
of them.
\writeup For $x>0$, $n^2xe^{-nx}\to0$ since the exponential beats $n^2$; at $x=0$ every term is $0$. With $u=nx$,
$\int_0^\infty n^2xe^{-nx}\dd x=\int_0^\infty ue^{-u}\dd u=1$. So $\int f_n=1\not\to0=\int\lim f_n$. The convergence is not uniform near $0$
($f_n(\frac1n)=\frac ne$), and there is no integrable bound: $\sup_nf_n(x)\ge f_n(x)$ with $n\approx\frac1x$ gives about $\frac1{ex}$, which
is not integrable at $0$. So the domination hypothesis fails.
\end{solution}

\begin{solution}{ex:gausscos}
\idea Differentiating brings down $-2x\sin(2ax)$, and $-2xe^{-x^2}$ is the derivative of $e^{-x^2}$: one integration by parts returns
the original integral.
\writeup Let $F(a)=\int_{\R}e^{-x^2}\cos(2ax)\dd x$. The integrand and its $a$-derivative $-2xe^{-x^2}\sin(2ax)$ are bounded by $e^{-x^2}$
and $2|x|e^{-x^2}$, both integrable, so by Theorem~\ref{thm:domination}
\[
\begin{aligned}F'(a)&=\int_{\R}(-2xe^{-x^2})\sin(2ax)\dd x\eqbr=\Bigl[e^{-x^2}\sin2ax\Bigr]_{-\infty}^\infty-2a\int_{\R}e^{-x^2}\cos2ax\dd x\eqbr=-2aF(a).\end{aligned}
\]
So $(e^{a^2}F)'=0$ and $F(a)=F(0)e^{-a^2}=\sqrt\pi e^{-a^2}$, using Example~\ref{ex:gaussfeynman}. In the language of Chapter~8: the
Gaussian is its own Fourier transform.
\end{solution}

\begin{solution}{ex:sin2}
\idea Integrate by parts so as to differentiate $\sin^2x$, which brings back a sine over $x$.
\writeup On $[\varepsilon,R]$, $\int_\varepsilon^R\frac{\sin^2x}{x^2}\dd x=\Bigl[-\frac{\sin^2x}x\Bigr]_\varepsilon^R+\int_\varepsilon^R\frac{\sin2x}x\dd x$. The bracket tends to
$0$ as $\varepsilon\to0$ and $R\to\infty$, and with $u=2x$ the last integral tends to $\int_0^\infty\frac{\sin u}u\dd u=\frac\pi2$. (This integral is
absolutely convergent, unlike Dirichlet's.)
\end{solution}

\begin{solution}{ex:xsinx}
\idea The factor $x$ is the only thing that breaks the symmetry $x\mapsto\pi-x$; averaging the integral with its reflection
replaces $x$ by the constant $\frac\pi2$.
\writeup Let $I$ be the integral. Since $\sin(\pi-x)=\sin x$ and $\cos^2(\pi-x)=\cos^2x$, the substitution gives
$I=\int_0^\pi\frac{(\pi-x)\sin x}{1+\cos^2x}\dd x$. Adding, $2I=\pi\int_0^\pi\frac{\sin x}{1+\cos^2x}\dd x$, and with $u=\cos x$ this is
$\pi\int_{-1}^1\frac{\mathrm du}{1+u^2}=\pi\cdot\frac\pi2$. So $I=\frac{\pi^2}4$.
\end{solution}

\begin{solution}{ex:loglog}
\idea Expand the logarithm that is regular at $0$ into its power series, integrate term by term against $\ln x$, and recognise
$\zeta(2)$ in what is left after partial fractions.
\writeup On $]0,1[$,
\[
\begin{aligned}\ln x\ln(1-x)&=(-\ln x)(-\ln(1-x))\eqbr=\sum_{n\ge1}\frac{x^n(-\ln x)}n,\end{aligned}
\]
a series of non-negative continuous functions. Its partial sums $S_N$ increase to the sum $S$, uniformly on every $[u,v]\subseteq\,]0,1[$ (the tail is at most
$|\ln u|\sum_{n>N}\frac{v^n}n$), and $0\le S_N\le S$ with $S$ continuous and integrable on $]0,1[$ (it is $O(x|\ln x|)$ at $0$ and $O(|\ln(1-x)|)$
at $1$). So Corollary~\ref{cor:dct} applies, and the integral is $\sum_n\frac1n\int_0^1x^n(-\ln x)\dd x=\sum_n\frac1{n(n+1)^2}$, using $\int_0^1x^n(-\ln x)\dd x=\frac1{(n+1)^2}$ (by parts).
With $\frac1{n(n+1)^2}=\frac1n-\frac1{n+1}-\frac1{(n+1)^2}$, the sum telescopes to $1-(\zeta(2)-1)=2-\frac{\pi^2}6$.
\end{solution}

\begin{solution}{ex:arctanparam}
\idea Differentiating in $a$ removes the logarithm and leaves a rational function of $x$, integrable by partial fractions.
\writeup Fix $A>0$. For $0\le a\le A$ the integrand is continuous, non-negative and at most $\frac{\ln(1+A^2x^2)}{1+x^2}$, which is
integrable (it is $O(\frac{\ln x}{x^2})$ at infinity); so $J$ is continuous on $[0,A]$ by Theorem~\ref{thm:domination}, and $J(0)=0$.
Its $a$-derivative is $\frac{2ax^2}{(1+a^2x^2)(1+x^2)}$. Near $a=0$ there is no good bound, so take $0<a_0\le a\le A$: then
$1+a^2x^2\ge a_0^2x^2$ and the derivative is at most $\frac{2A}{a_0^2(1+x^2)}$, integrable. So $J$ is $C^1$ on $]0,\infty[$ with
$J'(a)=\int_0^\infty g_a(x)\dd x$, where $g_a(x)=\frac{2ax^2}{(1+a^2x^2)(1+x^2)}$. For $a\ne1$,
\[
g_a(x)=\frac{2a}{a^2-1}\Bigl(\frac1{1+x^2}-\frac1{1+a^2x^2}\Bigr),
\]
so $J'(a)=\frac{2a}{a^2-1}\bigl(\frac\pi2-\frac\pi{2a}\bigr)=\frac\pi{1+a}$, and by continuity of $J'$ also at $a=1$. So $J(a)=\pi\ln(1+a)+C$ on
$]0,\infty[$; letting $a\to0$ and using the continuity of $J$ at $0$ gives $C=0$, so $J(a)=\pi\ln(1+a)$.
\end{solution}

\section*{Sources and further reading}

\begin{itemize}
\item R.~Feynman, \emph{Surely You're Joking, Mr.~Feynman!}, Norton, 1985, the chapter ``A Different Box of Tools''.
\item P.~Nahin, \emph{Inside Interesting Integrals}, 2nd ed., Springer, 2020. Hundreds of integrals done with Feynman's trick,
symmetry and Frullani.
\item G.~Boros and V.~Moll, \emph{Irresistible Integrals}, Cambridge University Press, 2004.
\item W.~Rudin, \emph{Principles of Mathematical Analysis}, 3rd ed., McGraw-Hill, 1976, Chapters~6 and~7 (integrals depending
on a parameter, uniform convergence).
\item J.~Arias de Reyna, \emph{On the theorem of Frullani}, Proceedings of the American Mathematical Society 109 (1990),
165--175.
\item D.~Bressoud, \emph{A Radical Approach to Lebesgue's Theory of Integration}, Cambridge University Press, 2008, for the
history of interchanging limits and integrals.
\end{itemize}
